\begin{definition}\label{conser} 
Un campo vectorial continuo $\mathbf{F}: A \subseteq \mathbb{R}^n \to \mathbb{R}^n$, con $A$ abierto y arcoconexo, se llama \emph{campo conservativo} si existe una función escalar $f: A \subseteq \mathbb{R}^n \to \mathbb{R}$ de clase $\mathcal{C}^1$ tal que:
\[ \nabla f(\mathbf{x}) = \mathbf{F}(\mathbf{x}), \quad \forall \mathbf{x} \in A \] 
En tal caso, el campo escalar $f$ se llama \emph{función potencial} del campo vectorial $\mathbf{F}$. 
\end{definition}

\begin{propertie}\label{conserv}  Sea $A\subset \Rn{n}$   ($n=2$ o $n=3$)  abierto y arcoconexo y sea  $\mathbf{F}:A  \to \Rn{n}$  de clase $C^{1}(A).$  Si $\mathbf{F}$ es  un campo conservativo entonces $$\nabla \times \mathbf{F}(\mathbf{x}) =0 \:\:\: \forall \mathbf{x} \in A.  $$  
\end{propertie}

\begin{example} Determinar si $\mathbf{F}(x,y,z) = (z\cos xz,\; z,\; y + x\cos xz)$ admite función potencial y, cn caso afirmativo,  hallarla.

Un campo vectorial $\mathbf{F} = (P, Q, R)$
de clase $\mathcal{C}^1$ en un abierto simplemente conexo admite función potencial
si y sólo si $\nabla \times \mathbf{F} = \mathbf{0}$, es decir:
\[
    \frac{\partial R}{\partial y} = \frac{\partial Q}{\partial z}, \qquad
    \frac{\partial P}{\partial z} = \frac{\partial R}{\partial x}, \qquad
    \frac{\partial Q}{\partial x} = \frac{\partial P}{\partial y}.
\]
con $P = z\cos xz$, $Q = z$, $R = y + x\cos xz$.

\[
    \frac{\partial R}{\partial y} = 1, \qquad
    \frac{\partial Q}{\partial z} = 1
\]
 
\[
    \frac{\partial P}{\partial z} = \cos xz + z \cdot (-x\sin xz) = \cos xz - xz\sin xz,
\]

\[
    \frac{\partial R}{\partial x} = \cos xz + x \cdot (-z\sin xz) = \cos xz - xz\sin xz
\]

\[
    \frac{\partial Q}{\partial x} = 0, \qquad
    \frac{\partial P}{\partial y} = 0
\]

Las tres condiciones se satisfacen, luego $\mathbf{F}$ \textbf{admite función potencial}.

\medskip
\textbf{Hallamos $\varphi$ tal que $\nabla\varphi = \mathbf{F}$.} Es decir:
\[
    \frac{\partial \varphi}{\partial x} = z\cos xz, \qquad
    \frac{\partial \varphi}{\partial y} = z, \qquad
    \frac{\partial \varphi}{\partial z} = y + x\cos xz.
\]

\[
    \varphi(x,y,z)
    = \int z\cos xz\, dx
    = z \cdot \frac{\sin xz}{z} + g(y,z)
    = \sin xz + g(y,z).
\]

\[
    \frac{\partial \varphi}{\partial y} = \frac{\partial g}{\partial y} = z
    \implies g(y,z) = yz + h(z).
\]
Luego:
\[
    \varphi(x,y,z) = \sin xz + yz + h(z).
\]

\[
    \frac{\partial \varphi}{\partial z} = x\cos xz + y + h'(z) = y + x\cos xz
    \implies h'(z) = 0 \implies h(z) = C.
\]

\textbf{Función potencial:}
\[
    \boxed{\varphi(x,y,z) = \sin xz + yz + C.}
\]
\end{example}

\begin{example} 
Determinar si el campo vectorial $\mathbf{F}(x,y,z) = (z\cos xz,\; z,\; y + x\cos xz)$ admite función potencial y, en caso afirmativo, hallarla.
\end{example}

\noindent\textbf{Solución:}  El campo vectorial $\mathbf{F} = (P, Q, R)$ está definido en todo $\mathbb{R}^3$ y sus componentes son combinaciones de funciones polinómicas y trigonométricas de clase $\mathcal{C}^1$. 
De acuerdo con la Propiedad~\ref{conserv}, si un campo es conservativo, sus derivadas parciales cruzadas deben ser iguales (lo que equivale a que su rotacional sea nulo).  Por lo tanto, antes de buscar la función potencial, verificamos si se cumplen estas condiciones necesarias:
\[ \frac{\partial R}{\partial y} = \frac{\partial Q}{\partial z}, \qquad \frac{\partial P}{\partial z} = \frac{\partial R}{\partial x}, \qquad \frac{\partial Q}{\partial x} = \frac{\partial P}{\partial y} \]

Identificamos las componentes del campo:
\[ P(x,y,z) = z\cos xz, \qquad Q(x,y,z) = z, \qquad R(x,y,z) = y + x\cos xz \]

Calculamos las derivadas parciales correspondientes:
\begin{align*}
    \frac{\partial R}{\partial y} &= 1, \qquad \frac{\partial Q}{\partial z} = 1 \implies \frac{\partial R}{\partial y} = \frac{\partial Q}{\partial z} \\
    \frac{\partial P}{\partial z} &= \cos xz - xz\sin xz, \qquad \frac{\partial R}{\partial x} = \cos xz - xz\sin xz \implies \frac{\partial P}{\partial z} = \frac{\partial R}{\partial x} \\
    \frac{\partial Q}{\partial x} &= 0, \qquad \frac{\partial P}{\partial y} = 0 \implies \frac{\partial Q}{\partial x} = \frac{\partial P}{\partial y}
\end{align*}

Como las condiciones necesarias se satisfacen, el campo $\mathbf{F}$ es un candidato válido para ser conservativo. Procederemos a intentar construir su función potencial de forma directa en el siguiente paso. Si logramos hallarla con éxito, quedará demostrado que el campo efectivamente admite función potencial.



Buscamos una función escalar $\varphi(x,y,z)$ tal que $\nabla\varphi = \mathbf{F}$. Esto nos plantea el siguiente sistema de ecuaciones diferenciales parciales:
\[ \text{(I)} \ \frac{\partial \varphi}{\partial x} = z\cos xz, \qquad \text{(II)} \ \frac{\partial \varphi}{\partial y} = z, \qquad \text{(III)} \ \frac{\partial \varphi}{\partial z} = y + x\cos xz \]

Comenzamos integrando la ecuación (I) respecto a la variable $x$:
\[ \varphi(x,y,z) = \int z\cos xz\, dx =\sin xz + g(y,z) \]
donde $g(y,z)$ representa la "constante" de integración, la cual puede depender de las variables que se mantuvieron fijas ($y$ y $z$).

Para determinar $g(y,z)$, derivamos la expresión obtenida respecto a $y$ e igualamos con la ecuación (II):
\[ \frac{\partial \varphi}{\partial y} = \frac{\partial g}{\partial y} = z \]
Integrando ahora respecto a $y$, hallamos la forma de la función $g$:
\[ g(y,z) = \int z\, dy = yz + h(z) \]
donde $h(z)$ es una nueva función que depende únicamente de $z$. Sustituyendo esto en nuestra expresión para $\varphi$, tenemos:
\[ \varphi(x,y,z) = \sin xz + yz + h(z) \]

Finalmente, para hallar $h(z)$, derivamos respecto a $z$ e igualamos con la ecuación (III):
\[ \frac{\partial \varphi}{\partial z} = x\cos xz + y + h'(z) = y + x\cos xz \]
Simplificando los términos comunes de ambos lados, obtenemos:
\[ h'(z) = 0 \implies h(z) = C, \quad \text{con } C \in \mathbb{R} \]

\noindent \textbf{Resultado:} La familia de funciones potenciales para el campo vectorial $\mathbf{F}$ es:
\[ \boxed{\varphi(x,y,z) = \sin xz + yz + C} \]



